% Generated draft: requires mathematical and source review.
\documentclass[11pt]{article}
\usepackage[margin=25mm]{geometry}
\usepackage{fontspec}
\setmainfont{Latin Modern Roman}
\usepackage{amsmath,amssymb,mathtools,unicode-math}
\setmathfont{Latin Modern Math}
\usepackage{enumitem,needspace,xurl,hyperref}
\setlength{\parindent}{0pt}
\setlength{\parskip}{5pt}
\setlength{\emergencystretch}{3em}
\newcommand{\N}{\mathbb{N}}
\newcommand{\Z}{\mathbb{Z}}
\newcommand{\Q}{\mathbb{Q}}
\newcommand{\R}{\mathbb{R}}
\newcommand{\C}{\mathbb{C}}
\newcommand{\F}{\mathbb{F}}
\newcommand{\A}{\mathbb{A}}
\newcommand{\PP}{\mathbb{P}}
\newcommand{\E}{\mathbb{E}}
\newcommand{\eps}{\varepsilon}
\newcommand{\dd}{\,\mathrm{d}}
\begin{document}
\section*{UnsolvedMath ID: 30004684}
\textbf{Source problem code:} catalogue entry
\subsection{Definitions and mathematical statement}
Let \(X\) denote the unspecified state space of the dynamical system in the supplied record. A collective variable is an observable map \(\xi:X\to Y\), where \(Y\) is an intended lower-dimensional space. If dimensions \(d=\dim X\) and \(k=\dim Y\) are defined, the intended dimensional reduction requires \(k<d\); the source gives no quantitative meaning to ``low-dimensional.'' These symbols describe the requested construction without imposing a particular class of state spaces or dynamics.

\textbf{Research question.} Given only a high-dimensional description of a dynamical system, how can one construct a low-dimensional collective variable that reproduces its dominant dynamical timescales?

\textbf{Unresolved formulation.} The source does not specify the admissible dynamical systems, the information contained in their descriptions, a definition of dominant dynamical timescales, or how timescales are to be assigned to the collective variable. It also does not specify whether reproduction means exact equality or approximation, an error criterion, or a reduced evolution law. Consequently, the question is a research objective rather than a fully quantified existence or algorithmic statement. No universal existence of such a collective variable is asserted.
\subsection{Short English statement}
Construct a low-dimensional observable from a high-dimensional system description so that it retains the system's dominant dynamical timescales. The allowed systems, input information, meaning of timescale preservation, and accuracy requirements remain unspecified. The supplied label is ``partially solved''; this draft does not independently verify that assessment.
\subsection{Sources}
\begin{itemize}
\item Local registry: \texttt{data/unsolvedmath/problems.json}, permanent UnsolvedMath ID \(30004684\), source problem code \texttt{catalogue entry}, title ``Reduced Variables Preserving Dominant Dynamical Timescales,'' in the collection ``Oberwolfach Reports Open Problems.'' The record assigns proposed year 2021 and supplies no named proposer. Its statement-status field is ``exact,'' with a supplied assertion of byte-for-byte retention; neither the transcription nor the external provenance has been independently verified here.
\item Supplied report citation: \emph{Mini-Workshop: Analysis of Data-driven Optimal Control (2022)}, Oberwolfach Reports, DOI 10.4171/owr/2021/23, pp.~1209--1257. Also listed as ``Oberwolfach Report 23/2021.'' \url{https://doi.org/10.4171/owr/2021/23}. The background describes Oberwolfach Reports as documenting mathematical workshops and their open-problem discussions. The supplied citation's 2022 designation and its 2021 report identifier require reconciliation.
\item Supplied literature reference: ``Husic and Pande, collective variables.'' \url{https://doi.org/10.1016/j.sbi.2017.02.006}.
\item Supplied literature reference: ``Rydzewski et al., reweighted manifold learning.'' \url{https://pubs.acs.org/doi/10.1021/acs.jctc.2c00873}.
\item Supplied literature assessment, checked on exactly \texttt{2026-08-22}: status ``partially solved.'' The supplied curation states that constructive transition-density and manifold-learning procedures, variational slow-coordinate methods, and reweighted data-driven approaches find collective variables under structural assumptions, while no general method works from arbitrary high-dimensional descriptions. This is supplied dated curation derived from the OWR material, not an independently verified literature assessment. Its time-sensitive status and cited claims require independent verification.
\end{itemize}
\end{document}
